% RetroNNetScaler --- IEEEtran conference paper
% Build:
%   pdflatex retronnetscaler.tex
%   bibtex retronnetscaler
%   pdflatex retronnetscaler.tex
%   pdflatex retronnetscaler.tex
\documentclass[conference]{IEEEtran}

\IEEEoverridecommandlockouts
\usepackage{cite}
\usepackage{amsmath,amssymb,amsfonts}
\usepackage{mathtools}
\usepackage{algorithmic}
\usepackage{algorithm}
\usepackage{graphicx}
\usepackage{textcomp}
\usepackage{xcolor}
\usepackage{booktabs}
\usepackage{listings}
\usepackage{url}
\usepackage{hyperref}

\hyphenation{op-tical net-works semi-conduc-tor}

\lstset{
  basicstyle=\ttfamily\scriptsize,
  breaklines=true,
  columns=fullflexible,
  keepspaces=true,
  aboveskip=0.4em,
  belowskip=0.4em,
  frame=single,
  rulecolor=\color{black!30},
}

\begin{document}

\title{RetroNNetScaler: Contour Upscaling for Hard-Palette NES Pixel Art}

\author{\IEEEauthorblockN{RetroNNetScaler Project}
\IEEEauthorblockA{\textit{Technical Report}\\
\texttt{https://github.com/amirsdream/RetroNNetScaler}\\
August 2026}
}

\maketitle

\begin{abstract}
Upscaling palette-limited Nintendo Entertainment System (NES) frames for modern
4K displays differs fundamentally from natural-image super-resolution.
Continuous filters and photograph-oriented networks invent mid-tone ``cream''
colours and silky outlines that erase the aliased aesthetic of 8-bit games.
We present \textbf{RetroNNetScaler}, a Windows realtime pipeline that embeds a
NES core (SimpleNES), presents frames through Direct3D~11, and offers three
hard-palette production modes: \emph{Clarity} (maximum integer nearest-neighbour),
\emph{Smooth} (Hyllian-style xBR-sharp tip wedges with neighbourhood colour snap),
and \emph{Contour} (Kopf-style planarized diagonals with \emph{shared global
half-planes} so that long $45^\circ$ stairs become a single continuous edge).
We further document a TinySR neighbour-classification CNN ($\approx\!52$k
weights) trained against Scale2xSFX teachers: palette drift can be forced to
zero by construction, yet unconstrained decoding produces multicolour edge
fringes. The central methodological result is that one should
\textbf{discretise the output domain first}; for continuous diagonals without
cream, adjacent cells must agree on one geometric edge equation rather than
regress RGB. Contour requires no learned weights at runtime and maps naturally
to an FPGA streaming datapath.
\end{abstract}

\begin{IEEEkeywords}
pixel-art upscaling, NES, Contour, Kopf--Lischinski, xBR, hard palette,
realtime D3D11, discrete super-resolution, FPGA.
\end{IEEEkeywords}

% ---------------------------------------------------------------------------
\section{Introduction}
\label{sec:intro}

Classic console graphics were authored on a square pixel lattice with a small
discrete palette. When those frames are shown on a $3840{\times}2160$ panel,
naive bilinear or bicubic filters blur edges and invent colours that never
existed in the ROM. Deep super-resolution models trained for photographs or
anime with continuous losses tend to produce a \emph{silky} look: soft
outlines, chroma bleed, and palette drift. For retro enthusiasts and
preservationists, that look is a regression rather than an improvement.

Rule-based pixel-art scalers such as EPX/Scale2x, Scale2xSFX, HQX, and
xBR/xBRZ address diagonals by \emph{copying} existing neighbour colours rather
than interpolating~\cite{mazzoleni,sfx,hyllian}. Kopf and Lischinski
depixelization goes further toward resolution-independent
vectors~\cite{kopf2011}, which many players perceive as too smooth for
classic NES draw. Our product goal is narrower: hard colours (no mid-tone
cream); cleaner, continuous $45^\circ$ diagonals; realtime 4K presentation;
and a path toward FPGA implementation.

This paper documents the RetroNNetScaler prototype: system architecture
(Section~\ref{sec:arch}), Clarity / Smooth / Contour algorithms
(Section~\ref{sec:methods}), the shared half-plane insight that fixes EPX
zigzags (Section~\ref{sec:epx}), implementation notes
(Section~\ref{sec:impl}), the TinySR CNN negative result
(Section~\ref{sec:neural}), discussion (Section~\ref{sec:discuss}), and an
FPGA-oriented roadmap (Section~\ref{sec:fpga}).

% ---------------------------------------------------------------------------
\section{Related Work}
\label{sec:related}

\subsection{Classical Pixel-Art Scalers}
EPX/Scale2x expands each pixel into a $2{\times}2$ block using equality tests
on the $3{\times}3$ neighbourhood and never invents colours~\cite{mazzoleni}.
Scale2xSFX (and Scale3xSFX) refine diagonal decisions with wider support and
reduce hollow-diamond artefacts~\cite{sfx}. xBR and xBRZ detect edges with
weighted colour distances (often in YUV) and write tip wedges; Hyllian-style
sharp variants avoid blending~\cite{hyllian}. These methods remain the gold
standard for a hard-palette look.

\subsection{Depixelization and Continuous SR}
Kopf and Lischinski~\cite{kopf2011} build a similarity graph, planarize
crossing diagonals with curves / sparse / islands heuristics, reshape cells
via a generalized Voronoi diagram on nodes and half-edges, then fit contour
splines. Kreuzer \emph{et~al.}~\cite{kreuzer2015} target realtime GPU
paths. Full spline fitting yields the classic ``inked'' vector look; we adopt
the \emph{planarization + hard ownership} stage and deliberately omit soft
B-spline shading so fills stay flat NES colours.

Deep unsupervised pixelization and modern GAN/transformer SR (e.g., Real-ESRGAN,
SwinIR) optimise perceptual or continuous losses and routinely violate discrete
palette integrity. ESPCN-style networks~\cite{shi2016} are attractive for
FPGA targets because of small parameter counts; we adopt that footprint in
Section~\ref{sec:neural} but replace RGB residual regression with a 9-way
neighbour classifier.

% ---------------------------------------------------------------------------
\section{System Architecture}
\label{sec:arch}

The PC prototype is organised as a single process:
\begin{lstlisting}
NES core (SimpleNES)
     |  256 x 240 RGBA @ ~60.098 Hz
     v
D3D11 upload -> scaler CS -> integer letterbox present
     |
Modes: Clarity | Smooth | Contour | NIS | Nearest | Bilinear
\end{lstlisting}

A custom Win32 launcher provides ROM selection (browse / drag-and-drop), mode
chips (Clarity / Smooth / Contour), and output resolution ($1080$p / $1440$p /
$4$K). Preferences persist under AppData. Gameplay is paced to the NES NTSC
rate independent of panel refresh. Presentation uses the maximum integer scale
factor that fits the backbuffer, with phase-locked letterboxing (OSSC-style),
so output pixels remain uniform width~\cite{bestupscale}.

\begin{table}[t]
  \caption{Scaling modes in the current build.}
  \label{tab:modes}
  \centering
  \renewcommand{\arraystretch}{1.15}
  \begin{tabular}{@{}p{1.35cm}p{3.1cm}p{2.7cm}@{}}
    \toprule
    \textbf{Mode} & \textbf{Algorithm} & \textbf{Palette} \\
    \midrule
    Clarity & Max integer nearest & Exact LQ \\
    Smooth  & xBR-sharp + snap & Hard 3$\times$3 \\
    Contour & Shared half-planes & Exact neighbour \\
    NIS     & NVIDIA Image Scaling & Continuous \\
    Nearest / Bilinear & Baselines & Exact / blend \\
    \bottomrule
  \end{tabular}
\end{table}

% ---------------------------------------------------------------------------
\section{Methods}
\label{sec:methods}

\subsection{Clarity: Integer Oracle}
\label{sec:clarity}
Clarity selects
\begin{equation}
  s = \min\bigl(\lfloor W/w_{\mathrm{src}}\rfloor,\;
                 \lfloor H/h_{\mathrm{src}}\rfloor\bigr),
\end{equation}
point-samples NES texels into a canvas of size
$(w_{\mathrm{src}}\,s)\times(h_{\mathrm{src}}\,s)$, and letterboxes into the
backbuffer. At $1080$p, NES $256{\times}240$ typically yields $s{=}4$
($1024{\times}960$); at $4$K, $s{=}9$ ($2304{\times}2160$). Clarity is the
colour-perfect reference: blocky, but any subsequent filter must improve
diagonals without introducing cream.

\subsection{Smooth: Sharp xBR with Hard Snap}
\label{sec:smooth}
Smooth implements a Hyllian-inspired xBR-sharp pass in HLSL. For each output
sample we evaluate edge dominance with a YUV colour distance, restrict writes
to outer tip wedges of each quadrant, apply an outline guard on high-contrast
corners, and finally apply \texttt{snapExact9}: the result is replaced by the
nearest of the nine low-resolution (LQ) neighbourhood colours
$\{A,\ldots,I\}$. Mid-tones cannot survive the snap, which eliminates cream
while keeping intentional tip geometry.

\subsection{Contour: Shared Half-Plane Diagonals}
\label{sec:contour}
Contour is the production answer to ``sharp fills, continuous edges.'' For each
output sample inside LQ cell $E$ at integer coordinates $(s_x,s_y)$ with
fractional offsets $(f_x,f_y)\in[0,1)^2$, let
\begin{equation}
  \mathbf{p} = (s_x+f_x,\; s_y+f_y),\quad
  U = p_x + p_y,\quad
  V = p_x - p_y.
\end{equation}
The algorithm proceeds as follows:
\begin{enumerate}
  \item Fetch the $3{\times}3$ neighbourhood $A{\ldots}I$ and test similarity
        with Kopf/hqx YUV thresholds~\cite{kopf2011}:
        $|\Delta Y|\le 48/255$, $|\Delta U|\le 7/255$, $|\Delta V|\le 6/255$.
  \item Planarize each of the four $2{\times}2$ blocks around $E$
        (Algorithm~\ref{alg:planar}): keep at most one diagonal
        ($A$--$E$ or $B$--$D$, etc.), using a local curves / sparse / islands
        score when both diagonals cross.
  \item Assign colour by a \textbf{global} half-plane on $U$ (slash $/$) or
        $V$ (backslash $\backslash$), with bisectors at half-integers so the
        dark cell and the light cell of a stair share the \emph{same} edge
        equation (Section~\ref{sec:epx}).
  \item If no $45^\circ$ rule fires, apply $2{:}1$ / $1{:}2$ hard wedges
        (still exact neighbour colours).
\end{enumerate}
The output colour is always one of $A{\ldots}I$---never a blend. No network
weights are required at runtime.

\begin{algorithm}[t]
\caption{Planar diagonal selection in a $2{\times}2$ block.}
\label{alg:planar}
\begin{algorithmic}[1]
\REQUIRE colours $a,b,d,e$ of
$\bigl[\begin{smallmatrix}a&b\\d&e\end{smallmatrix}\bigr]$
\ENSURE $0$ (none), $1$ (keep $a$--$e$), or $2$ (keep $b$--$d$)
\STATE $ae \leftarrow \mathrm{Similar}(a,e)$;\;
       $bd \leftarrow \mathrm{Similar}(b,d)$
\IF{$\neg ae \wedge \neg bd$}
  \RETURN $0$
\ENDIF
\IF{all four cardinal edges similar}
  \RETURN $0$ \COMMENT{fully connected shade}
\ENDIF
\IF{$ae \wedge \neg bd$}
  \RETURN $1$
\ENDIF
\IF{$bd \wedge \neg ae$}
  \RETURN $2$
\ENDIF
\IF{any cardinal similar}
  \RETURN $0$ \COMMENT{drop crossing diagonals}
\ENDIF
\STATE score both diagonals (curves / sparse / islands)
\RETURN higher-scoring diagonal, or $0$ on tie
\end{algorithmic}
\end{algorithm}

% ---------------------------------------------------------------------------
\section{Why Local EPX Cuts Still Look Pixelated}
\label{sec:epx}

Classic EPX paints a triangle inside each cell, e.g.,
$f_x+f_y < \tfrac12$ for a top-left cut. On a long slash stair, the dark
cell's cut and the light cell's cut often sit on \textbf{parallel but offset}
lines---different constants of $U=x+y$. The visible border becomes a thick
zigzag at the scale of \emph{source} pixels: still ``pixelated,'' even though
each cut is continuous inside its cell.

Contour forces both sides onto one shared bisector. With
$U_s = s_x+s_y$ and $V_s = s_x-s_y$:
\begin{align}
  \text{TL keep }B\text{--}D:&\quad
    U < U_s + \tfrac12 \;\Rightarrow\; B,\\
  \text{BR keep }F\text{--}H:&\quad
    U > U_s + \tfrac32 \;\Rightarrow\; F,\\
  \text{TR keep }B\text{--}F:&\quad
    V > V_s + \tfrac12 \;\Rightarrow\; B,\\
  \text{BL keep }D\text{--}H:&\quad
    V < V_s - \tfrac12 \;\Rightarrow\; D.
\end{align}
Complementary corners of a stair therefore agree on a single straight
$45^\circ$ interface (e.g., $U = U_s+\tfrac12$ meeting
$U = U_s+\tfrac32$ from the adjacent cell), recovering continuity without
blending.

% ---------------------------------------------------------------------------
\section{Implementation Notes}
\label{sec:impl}

Clarity, Smooth, and Contour share one compute shader
(\texttt{shaders/tiny\_sr.hlsl}), selected by a mode constant. Contour
dispatches at destination resolution with integer scale $s$ from the present
path. The shader is compiled at runtime with \texttt{D3DCompile}
(\texttt{cs\_5\_0}). The Win32 launcher maps the Contour chip to the internal
\texttt{ScaleMode::Neural} enum for preferences compatibility; the CLI accepts
\texttt{--scale neural}.

Integer phase lock follows the practice summarised in the project notes:
content size exactly $w_{\mathrm{src}}s \times h_{\mathrm{src}}s$,
integer-centred viewport, point sampling after the scaler, and optional
one-pixel phase nudges~\cite{bestupscale}.

% ---------------------------------------------------------------------------
\section{The Neural Experiment (Negative Result)}
\label{sec:neural}

Before Contour, we trained a TinySR network to predict, for each LQ pixel, a
class in $\{A,\ldots,I\}$ for every subpixel of a $4{\times}4$ block
(pixel-shuffle head). Decoded colour is exactly
$\mathrm{neighbour}[\arg\max]$---zero palette drift by construction.
\begin{lstlisting}
Conv3x3 C=32 ReLU -> Conv3x3 C=32 ReLU
  -> Conv3x3 9*4*4 -> pixel_shuffle(r=4)
  -> logits (9,H*4,W*4)   ~51.8k params
\end{lstlisting}
Labels came from a Scale2xSFX teacher chained to $4{\times}$. Non-centre
classes are rare ($\approx\!0.3$--$1\%$ of pixels), so training used
class-weighted cross-entropy with a large edge multiplier, edge-biased crops,
AdamW, and early stopping on edge-class accuracy (CUDA, GeForce RTX~4080).

\begin{table}[t]
  \caption{Classifier validation on Super Mario Bros.\ hold-out frames
  (illustrative means).}
  \label{tab:cnn}
  \centering
  \renewcommand{\arraystretch}{1.15}
  \begin{tabular}{@{}lccc@{}}
    \toprule
    \textbf{Metric} & \textbf{NN} & \textbf{CNN} & \textbf{SFX} \\
    \midrule
    Acc.\ (all)  & --- & $\approx 96\%$ & $100\%$ \\
    Acc.\ (edge) & --- & $\approx 72\%$ & $100\%$ \\
    Palette drift & $0$ & $0$ & $0$ \\
    \bottomrule
  \end{tabular}
\end{table}

\subsection{Failure Modes}
\paragraph{Cream.}
Residual RGB regression underfits sparse EPX/SFX deltas; raising gain without
discrete decoding invents mid-tones.

\paragraph{Fringes.}
At roughly $72\%$ edge accuracy, false non-$E$ classes copy a wrong palette
entry (e.g., HUD red into a black silhouette), appearing as multicolour static
at $4{\times}$. Tip gates and a runtime fallback to Smooth mitigated symptoms
but did not beat Contour's geometric consistency.

\paragraph{Lesson.}
For this medium, learn (or hard-code) \emph{geometry} rather than continuous
colour. Sparse teacher signal---only a fraction of a percent of HQ pixels
differ from nearest under EPX/SFX---makes unweighted continuous losses
dominated by flat regions.

% ---------------------------------------------------------------------------
\section{Discussion}
\label{sec:discuss}

Contour targets $45^\circ$ stairs and shallow $2{:}1$ wedges. Very irregular
silhouettes and text may still benefit from Smooth's tip logic or Clarity's
honest blocks. At modest integer scales ($s{=}4$ at $1080$p), even a
mathematically straight hard edge is a one-pixel staircase of \emph{output}
samples; that is rasterization, not the source-pixel zigzag Contour removes.
Higher $s$ (4K) makes continuous diagonals more visually obvious.

We intentionally omit Kopf spline optimisation: curved inked contours read as
silky to many NES players. Hard half-planes keep the look closer to ``drawn
with exact colours'' than ``vectorized illustration.''

% ---------------------------------------------------------------------------
\section{FPGA and Future Work}
\label{sec:fpga}

Contour maps to a streaming pipeline: line buffers for a $3{\times}3$ window,
YUV similarity comparators, four planar-diagonal units, and a few plane tests
per output sample---no DRAM weight traffic. Clarity is already trivial on
FPGA. Intended evaluation board: DE10-Nano for Clarity/Contour; larger devices
only if a learned tip classifier returns.

Future work includes: tightening Contour on shallow non-$45^\circ$ hills and
sprite silhouettes; an optional NES PAR ($8{:}7$) toggle for CRT-accurate
width; a tip-only learned gate distilled from Smooth if Contour gaps remain;
int8 / cycle-accurate RTL for the Contour datapath; and multi-game visual QA
beyond Super Mario Bros.

% ---------------------------------------------------------------------------
\section{Conclusion}
\label{sec:conclusion}

RetroNNetScaler is a complete path from NES emulation to 4K presentation with
explicit respect for discrete palettes. Clarity is the honest oracle; Smooth
improves tips without cream; Contour makes long diagonals share one geometric
edge instead of EPX's offset zigzags. A FPGA-sized neighbour CNN achieved zero
palette drift but fringed until abandoned for Contour at runtime. The project's
central lesson is: \textbf{discretise first}; for continuous edges without
silk, make both sides of the stair agree on the same half-plane.

\section*{Acknowledgment}
Training and validation dumps must use legally obtained ROMs; ROMs are not
redistributed with the project.

\bibliographystyle{IEEEtran}
\bibliography{references}

\end{document}
